\documentclass[11pt,a4paper]{article}
\usepackage[T1]{fontenc}
\usepackage[utf8]{inputenc}
\usepackage{amsmath,amssymb,amsthm}
\usepackage[margin=1in]{geometry}
\usepackage[colorlinks=true,linkcolor=blue,urlcolor=blue]{hyperref}
\theoremstyle{plain}
\newtheorem{theorem}{Theorem}
\newtheorem{lemma}[theorem]{Lemma}
\newtheorem{proposition}[theorem]{Proposition}
\newtheorem{corollary}[theorem]{Corollary}
\theoremstyle{definition}
\newtheorem{remark}[theorem]{Remark}
% A quoted conjecture can be a single hundred-term formula whose terms are thirty-digit
% numbers. TeX cannot hyphenate those, so without a little stretch every such quote runs off
% the page; the linked-a-file papers are all of that shape.
\setlength{\emergencystretch}{3em}

\title{The empirical recurrence for OEIS A234127, proved by transfer matrix}
\author{Adrian Perez Fontelles\\ \small Independent researcher}
\date{14 September 2026}
\begin{document}
\maketitle

\begin{abstract}
OEIS A234127 counts $(n+1)\times5$ arrays over $\{0,\dots,6\}$ subject to a condition imposed on
every $2\times2$ subblock, and carries an empirical recurrence of order $78$ contributed
by R.~H.~Hardin. It is true, and it is not an empirical matter at all. The condition is local
to each subblock, so the rows of an admissible array are the vertices of a finite digraph
on $S=16807$ vertices whose edges are the admissible consecutive pairs, and the count is a
number of walks: $a(n)=\mathbf{1}^{\!\top}M^{n}\mathbf{1}$. Writing $q$
for the characteristic polynomial of the conjectured recurrence, the recurrence holds for all
$n$ exactly when $\mathbf{1}^{\!\top}M^{j}q(M)\mathbf{1}=0$ for every $j\ge0$, and the
Cayley--Hamilton theorem bounds that to $j<S$. It is therefore a finite computation in exact
integer arithmetic, and it comes out zero.
\end{abstract}

\noindent\small 2020 Mathematics Subject Classification. 05A15, 05B45, 15B36.\normalsize

\section{The sequence and the conjecture}

OEIS A234127 is ``Number of (n+1) X (4+1) 0..6 arrays with every 2 X 2 subblock having the absolute values of all six edge and diagonal differences no larger than 1.''. It has offset $1$ and begins
\[
12457, 484731, 19007583, 753253097, 29939517797, 1191882104721, 47493714739641, 1893857438251263,\ \dots
\]
The entry carries the following comment:
\begin{quote}\small
Empirical recurrence of order 78 (see link above).

Empirical: a(n) = 169*a(n-1) - 11674*a(n-2) + 414890*a(n-3) - 7427237*a(n-4) + 32804023*a(n-5) + 1117149338*a(n-6) - 18009100812*a(n-7) + 2056648790*a(n-8) + 1985726671261*a(n-9) - 10511696375014*a(n-10) - 110096293381822*a(n-11) + 1040939428734419*a(n-12) + 3485490424099116*a(n-13) - 56652801677310328*a(n-14) - 56799400580096979*a(n-15) + 2094767581549699564*a(n-16) - 45854162284354546*a(n-17) - 57133795688753687919*a(n-18) + 24615958300958618005*a(n-19) + 1197797296168301630588*a(n-20) - 571497373046657692896*a(n-21) - 19691038114553818164130*a(n-22) + 5833724596687105926873*a(n-23) + 255373495178957856784293*a(n-24) + 4748915716239887418869*a(n-25) - 2602850392671694560455668*a(n-26) - 1050074358043666358573187*a(n-27) + 20612871309330499223554394*a(n-28) + 16995757962032073085334565*a(n-29) - 124464415192454926426121719*a(n-30) - 159642872431016413534464040*a(n-31) + 556056236426121809829347797*a(n-32) + 1016907936279532078238106654*a(n-33) - 1739655976286643374035203803*a(n-34) - 4590617841272351697315772557*a(n-35) + 3310114549172089594435474644*a(n-36) + 14882502736708444962657730120*a(n-37) - 1395696151786626090501703530*a(n-38) - 34701580086875870061708067414*a(n-39) - 12266856914690124624562882957*a(n-40) + 57981811773677043522771759712*a(n-41) + 42103956511155070937589853264*a(n-42) - 68752894437027003723220706361*a(n-43) - 76817535317303726154767307039*a(n-44) + 56244730182915488760099463203*a(n-45) + 94451365594997429270893799665*a(n-46) - 28857451375436408117507086277*a(n-47) - 84357068305411477588604951302*a(n-48) + 5189846224860976307508560235*a(n-49) + 56905592685975310174178962466*a(n-50) + 5100866950486402420875505830*a(n-51) - 29773489943557320834258638981*a(n-52) - 5431846624528928907005461239*a(n-53) + 12335518243528613981802869744*a(n-54) + 2841937057460547271244875185*a(n-55) - 4113075177372064339874634536*a(n-56) - 980259443271902542619309938*a(n-57) + 1114097956913771690972926604*a(n-58) + 233513334083502497698912306*a(n-59) - 244615048416093365794159780*a(n-60) - 36898360518952099165476604*a(n-61) + 42712318948265171961115392*a(n-62) + 3112806715400777335976520*a(n-63) - 5692377564731435436490128*a(n-64) + 73836107262759445978144*a(n-65) + 538911014404599693250368*a(n-66) - 54123589064306956078656*a(n-67) - 31731569640816448883328*a(n-68) + 6224297067875670876928*a(n-69) + 813679443272237946880*a(n-70) - 306029778210944672256*a(n-71) + 8386779639879844864*a(n-72) + 4889886358315524096*a(n-73) - 392190064653606912*a(n-74) - 30054443916754944*a(n-75) + 3637884673327104*a(n-76) + 58398096752640*a(n-77) - 10901363097600*a(n-78)
\end{quote}
As of the ``Last modified'' line on the live entry (Sep 26 2025 20:45:35, revision 12) the statement is
still recorded as empirical, and nothing on the entry records it as settled either way.

Written out, the claim is
\begin{equation}\label{eq:conj}
a(n)\;=\;169\,a(n-1) - 11674\,a(n-2) + 414890\,a(n-3) - 7427237\,a(n-4) + 32804023\,a(n-5) + 1117149338\,a(n-6) - 18009100812\,a(n-7) + 2056648790\,a(n-8) + 1985726671261\,a(n-9) - 10511696375014\,a(n-10) - 110096293381822\,a(n-11) + 1040939428734419\,a(n-12) + 3485490424099116\,a(n-13) - 56652801677310328\,a(n-14) - 56799400580096979\,a(n-15) + 2094767581549699564\,a(n-16) - 45854162284354546\,a(n-17) - 57133795688753687919\,a(n-18) + 24615958300958618005\,a(n-19) + 1197797296168301630588\,a(n-20) - 571497373046657692896\,a(n-21) - 19691038114553818164130\,a(n-22) + 5833724596687105926873\,a(n-23) + 255373495178957856784293\,a(n-24) + 4748915716239887418869\,a(n-25) - 2602850392671694560455668\,a(n-26) - 1050074358043666358573187\,a(n-27) + 20612871309330499223554394\,a(n-28) + 16995757962032073085334565\,a(n-29) - 124464415192454926426121719\,a(n-30) - 159642872431016413534464040\,a(n-31) + 556056236426121809829347797\,a(n-32) + 1016907936279532078238106654\,a(n-33) - 1739655976286643374035203803\,a(n-34) - 4590617841272351697315772557\,a(n-35) + 3310114549172089594435474644\,a(n-36) + 14882502736708444962657730120\,a(n-37) - 1395696151786626090501703530\,a(n-38) - 34701580086875870061708067414\,a(n-39) - 12266856914690124624562882957\,a(n-40) + 57981811773677043522771759712\,a(n-41) + 42103956511155070937589853264\,a(n-42) - 68752894437027003723220706361\,a(n-43) - 76817535317303726154767307039\,a(n-44) + 56244730182915488760099463203\,a(n-45) + 94451365594997429270893799665\,a(n-46) - 28857451375436408117507086277\,a(n-47) - 84357068305411477588604951302\,a(n-48) + 5189846224860976307508560235\,a(n-49) + 56905592685975310174178962466\,a(n-50) + 5100866950486402420875505830\,a(n-51) - 29773489943557320834258638981\,a(n-52) - 5431846624528928907005461239\,a(n-53) + 12335518243528613981802869744\,a(n-54) + 2841937057460547271244875185\,a(n-55) - 4113075177372064339874634536\,a(n-56) - 980259443271902542619309938\,a(n-57) + 1114097956913771690972926604\,a(n-58) + 233513334083502497698912306\,a(n-59) - 244615048416093365794159780\,a(n-60) - 36898360518952099165476604\,a(n-61) + 42712318948265171961115392\,a(n-62) + 3112806715400777335976520\,a(n-63) - 5692377564731435436490128\,a(n-64) + 73836107262759445978144\,a(n-65) + 538911014404599693250368\,a(n-66) - 54123589064306956078656\,a(n-67) - 31731569640816448883328\,a(n-68) + 6224297067875670876928\,a(n-69) + 813679443272237946880\,a(n-70) - 306029778210944672256\,a(n-71) + 8386779639879844864\,a(n-72) + 4889886358315524096\,a(n-73) - 392190064653606912\,a(n-74) - 30054443916754944\,a(n-75) + 3637884673327104\,a(n-76) + 58398096752640\,a(n-77) - 10901363097600\,a(n-78)
\end{equation}
for all $n>78$.

\section{The condition is local, so the rows form a digraph}

Write an array of the shape in question as its list of rows
$r^{(1)},r^{(2)},\dots$, each an element of
$V=\{0,1,\dots,6\}^{5}$, so $|V|=S=16807$. A $2\times2$ subblock of the array
occupies two consecutive rows and two consecutive columns; if the two rows are $r$
and $s$ (in that order) and the two columns are numbered $j,j+1$,
the subblock is
\[
\begin{pmatrix}\alpha&\beta\\\gamma&\delta\end{pmatrix}=\begin{pmatrix}r_j&r_{j+1}\\s_j&s_{j+1}\end{pmatrix}.
\]
The entry's requirement on that subblock is
\begin{equation}\label{eq:loc}
\max\bigl(|\beta-\alpha|,|\delta-\beta|,|\gamma-\delta|,|\alpha-\gamma|,|\delta-\alpha|,|\gamma-\beta|\bigr)\;\le\;1,
\end{equation}
a condition on the four entries alone.

For $r,s\in V$ declare $r\to s$ an edge of a digraph $G$ when, for every
$1\le j\le 5-1$,
\begin{itemize}
\item the four entries above satisfy \eqref{eq:loc};

\end{itemize}
Let $M\in\{0,1\}^{S\times S}$ be the adjacency matrix of $G$ and $\mathbf{1}$ the
all-ones vector.

\begin{lemma}\label{lem:walk}
The arrays counted by the entry with $L$ rows are exactly the walks
$r^{(1)}\to r^{(2)}\to\cdots\to r^{(L)}$ of length $L-1$ in $G$. Consequently
\[
a(n)\;=\;\mathbf{1}^{\!\top}M^{n}\mathbf{1} .
\]
\end{lemma}

\begin{proof}
An array with $L$ rows and $5$ columns is precisely a sequence
$r^{(1)},\dots,r^{(L)}$ of elements of $V$; conversely any such sequence is an array.
Every $2\times2$ subblock lies in two consecutive rows $r^{(t)},r^{(t+1)}$ and two
consecutive columns, so the array satisfies the entry's condition on all of its subblocks
if and only if each consecutive pair $\bigl(r^{(t)},r^{(t+1)}\bigr)$ satisfies it for
every column pair --- which is exactly the edge relation of $G$. The number of walks of
length $L-1$, summed over all starting and ending vertices, is
$\mathbf{1}^{\!\top}M^{L-1}\mathbf{1}$, since $(M^{k})_{r,s}$ counts the walks of
length $k$ from $r$ to $s$. The entry's index $n$ corresponds to $L=n+1$ rows.
\end{proof}

\section{The criterion}

Let $q(t)=t^{78}-\sum_i c_i\,t^{78-i}$ be the characteristic polynomial of
\eqref{eq:conj}, with the $c_i$ read off that equation.

\begin{lemma}\label{lem:crit}
For every $n$ with $n+1-1\ge78$,
\[
a(n)-\sum_i c_i\,a(n-i)\;=\;\mathbf{1}^{\!\top}M^{\,n+1-1-78}\,q(M)\,\mathbf{1} .
\]
Hence \eqref{eq:conj} holds from that point on if and only if
$\mathbf{1}^{\!\top}M^{j}q(M)\mathbf{1}=0$ for every $j\ge0$; and it is enough to check
$j=0,1,\dots,S-1$.
\end{lemma}

\begin{proof}
By Lemma~\ref{lem:walk}, $a(n-i)=\mathbf{1}^{\!\top}M^{\,n+1-1-i}\mathbf{1}$, so
\[
a(n)-\sum_i c_i a(n-i)
=\mathbf{1}^{\!\top}\Bigl(M^{m}-\sum_i c_i M^{m-i}\Bigr)\mathbf{1}
=\mathbf{1}^{\!\top}M^{\,m-78}\,q(M)\,\mathbf{1}
\]
with $m=n+1-1$, which is the stated identity; the equivalence follows by letting
$m-78=j$ run over $\ge0$. For the last clause put $w=q(M)\mathbf{1}$. By the
Cayley--Hamilton theorem $M^{S}$ is an integer combination of $I,M,\dots,M^{S-1}$, so
$\mathbf{1}^{\!\top}M^{j}w$ for $j\ge S$ is the corresponding combination of the earlier
values; if those all vanish, so do all the rest.
\end{proof}

\section{The computation}

\begin{theorem}
The recurrence \eqref{eq:conj} holds for every $n>78$, which is the statement of
Section~1.
\end{theorem}

\begin{proof}
The digraph was constructed from \eqref{eq:loc} by a depth-first scan across the $5$
columns: the edge condition is a conjunction of one constraint per consecutive column
pair, so the successors of a row $r$ are enumerated column by column without testing all
$S^2$ pairs. The vector $w=q(M)\mathbf{1}\in\mathbb{Z}^{S}$ was then formed by $78$
matrix--vector products in exact integer arithmetic, and $\mathbf{1}^{\!\top}M^{j}w$ was
evaluated for $j=0,1,\dots$, again exactly. Every one of those integers vanishes from the
index corresponding to $n=78+1$ onwards, and the run continued until $w$ itself was the
zero vector, which settles every larger $j$ at once. By Lemma~\ref{lem:crit} the recurrence
holds for every $n>78$.
\end{proof}

\begin{remark}
The argument gives more than the recurrence: it shows the sequence is $C$-finite with
characteristic polynomial dividing that of $M$, so it has a rational generating function
whose denominator divides $\det(I-xM)$. The conjectured recurrence is one consequence; the
minimal one is a divisor of $q$.
\end{remark}

\section{Verification}

Three checks, each independent of the algebra above.

First, the digraph was built directly from the entry's stated condition and the walk counts
$\mathbf{1}^{\!\top}M^{L-1}\mathbf{1}$ compared against the entry's DATA at
every one of the $15$ published indices; the values agree exactly. That is what ties
the digraph to the sequence: a model that reproduced only some of the published terms would
be the wrong model, and would have been rejected. In particular it is what rules out a
misreading of the entry's wording, since a different reading of the condition gives a
different digraph and different counts.

Second, the recurrence \eqref{eq:conj} was evaluated on those published terms in exact
integer arithmetic, with no matrices involved, and vanishes wherever it applies.

Third, the annihilation test of Lemma~\ref{lem:crit} was run in exact integer arithmetic
until the working vector was identically zero, not to a sampled prefix, and returned zero
throughout.

\begin{thebibliography}{9}
\bibitem{oeis} The OEIS Foundation, \emph{The On-Line Encyclopedia of Integer
Sequences}, \url{https://oeis.org}, 2026. Sequence A234127.
\bibitem{stanley} R.~P.~Stanley, \emph{Enumerative Combinatorics}, Volume 1, 2nd ed.,
Cambridge University Press, 2011. (Transfer-matrix method, Section 4.7.)
\bibitem{fs} P.~Flajolet and R.~Sedgewick, \emph{Analytic Combinatorics}, Cambridge
University Press, 2009.
\bibitem{horn} R.~A.~Horn and C.~R.~Johnson, \emph{Matrix Analysis}, 2nd ed., Cambridge
University Press, 2013.
\end{thebibliography}

\end{document}
